\documentclass[handout]{beamer}
\mode<presentation>
\usepackage{xcolor}
\usepackage[toc]{multitoc}
\usepackage{color}
%\usepackage{enumitem}
\usetheme[secheader]{Boadilla}
\usecolortheme{whale}
\setbeamertemplate{caption}[numbered]
\setbeamertemplate{bibliography item}[text]
\setbeamertemplate{navigation symbols}{}

\theoremstyle{definition}
\newtheorem{definisi}[theorem]{Definisi}
\theoremstyle{example}
\newtheorem{contoh}[theorem]{Contoh}

\renewcommand*{\multicolumntoc}{2}

%#make sure to change this part, since it is predefined
      %\defbeamertemplate*{footline}{infolines theme}
      \setbeamertemplate{footline}
        {
      \leavevmode%
      \hbox{%
      \begin{beamercolorbox}[wd=.333333\paperwidth,ht=2.25ex,dp=1ex,center]{author in head/foot}%
        \usebeamerfont{author in head/foot}\insertshortauthor%~~(\insertshortinstitute)
      \end{beamercolorbox}%
      \begin{beamercolorbox}[wd=.433333\paperwidth,ht=2.25ex,dp=1ex,center]{title in head/foot}%
        \usebeamerfont{title in head/foot}\insertshorttitle
      \end{beamercolorbox}%
      \begin{beamercolorbox}[wd=.233333\paperwidth,ht=2.25ex,dp=1ex,right]{date in head/foot}%
        \usebeamerfont{date in head/foot}\insertshortdate{}\hspace*{2em} \insertframenumber{} / \inserttotalframenumber\hspace*{2ex} 
      \end{beamercolorbox}}%
      \vskip0pt%
    }

%\AtBeginSection[] {
%    \begin{frame}<beamer>{Daftar Isi}
%    \frametitle{Daftar Isi}
		%\tableofcontents[currentsection,currentsubsection,sectionstyle=show/hide,subsectionstyle=show/show/hide,subsubsectionstyle=show/show/show/hide]
   % \end{frame}
%}


\subtitle[Logika Matematika]{Logika Matematika}
\title[Logika Preposisi (Kalkulus Preposisi) ]{Logika Preposisi (Kalkulus Preposisi)}
\author{Cecep Suwanda, S.Si.,M.Kom.}
\date{Juni 9, 2025}
\institute{Universitas Bale Bandung}

\begin{document}

\maketitle

%\begin{frame}
%\titlepage
%\end{frame}

\section*{Daftar Isi}
\begin{frame}[shrink]{Daftar Isi}
\frametitle{Daftar Isi}
\tableofcontents
\end{frame}


\section{Definisi Logika}
\begin{frame}
\frametitle{Definisi Logika}
    \begin{definisi}
   	  Ilmu yang mempelajari tentang cara berpikir yang logis/masuk akal 
    \end{definisi}
	\begin{definisi}
      Logika matematika adalah	ilmu yang digunakan untuk menentukan nilai kebenaran dari suatu pernyataan atau penarikan kesimpulan 
      berdasarkan aturan-aturan dasar yang berlaku.         
    \end{definisi}

\end{frame}

\section{Kalimat Tertutup}
\begin{frame}
\frametitle{Kalimat Tertutup}

 \begin{definisi}[Kalimat Tertutup/Pernyataan]
    Adalah suatu kalimat yang bernilai benar saja atau salah saja. Dengan kata lain, tidak sekaligus kedua-duanya. Pernyataan disebut juga \textbf{preposisi}. Yang dapat di golongkan sebagai pernyataan adalah kalimat-kalimat yang menerangkan sesuatu (disebut : \textit{kalimat deklaratif})
 \end{definisi}
 
\begin{contoh}	
	\begin{enumerate}
		\item Menara itu tinggi. 
		\item Jumlah hari ada 7.
		\item 3 adalah bilangan ganjil
		\item nilai x yang memenuhi 3x+1=7 adalah 2
		\item 5 kurang dari 3
	\end{enumerate}	 
\end{contoh}
\end{frame}

\section{Kalimat Terbuka}
\begin{frame}
	\frametitle{Kalimat Terbuka}

\begin{definisi}[Kalimat Terbuka]
	 Kalimat yang memuat peubah/variabel, sehingga belum dapat ditentukan nilai kebenarannya ( benar atau salah ). Tetapi apabila variabel diganti nilai tertentu akan menjadi suatu pernyataan.
\end{definisi}

\begin{contoh}
	\begin{enumerate}
		\item $y>5$	
		\item $x + 3 = 8$	
	\end{enumerate}
\end{contoh}

\end{frame}

\section{Nilai Kebenaran Suatu Pernyataan}
\begin{frame}[shrink]
\frametitle{Nilai Kebenaran Suatu Pernyataan}
  \begin{block}{Dasar Empiris}
  	Menentukan benar atau salah dari sebuah pernyataan berdasarkan fakta yang ada atau dijumpai dalam kehidupan sehari-hari.
  \end{block}
  
  \begin{contoh}
  	\begin{enumerate}
  		\item Ibukota jawa Timur adalah Surabaya.
  		\item Air adalah benda padat.
  	\end{enumerate}
  \end{contoh}
  
  
  \begin{block}{Dasar tak Empiris}
  	Menentukan benar atau salah dari sebuah pernyataan dengan memakai bukti atau perhitungan-perhitungan dalam matematika.
  \end{block}

  \begin{contoh}
  	\begin{enumerate}
  		\item Akar persamaan $3x-1=5$ adalah 2.
  		\item Jika $x > 1$, maka $x > 2$.
  			
  	\end{enumerate}
  \end{contoh}
\end{frame}

\section{Lambang Suatu Pernyataan}
\begin{frame}
\frametitle{Lambang Suatu Pernyataan}

\begin{block}{Lambang Suatu pernyataan}
 Suatu pernyataan dilambangkan dengan memakai huruf kecil, seperti \textit{a}, \textit{b}, \textit{c},...,\textit{p},\textit{q},\textit{r},... dan seterusnya.
\end{block}

\begin{contoh}
	Pernyataan "4 adalah bilangan genap" dapat dilambangkan dengan memakai huruf \textit{p}.\\
	Ditulis:\\
	\textit{p} :  4 adalah bilangan genap.
\end{contoh}
\end{frame}

\section{Lambang Nilai Kebenaran Suatu Pernyataan}
\begin{frame}
\frametitle{Lambang Nilai Kebenaran Suatu Pernyataan}
\begin{block}{Lambang Nilai Kebenaran Suatu Pernyataan}
Nilai kebenaran dilambangkan dengan memakai huruf Yunani $\tau$ (dibaca: \textit{tau})
\end{block}
  \begin{contoh}
  	\begin{enumerate}
	 \item $\tau(\textit{p}) = T$ dibaca “nilai kebenaran pernyataan \textit{p} adalah \textit{T}” atau “pernyataan \textit{p} mempunyai nilai kebenaran \textit{T}”.
	\item $\textit{q}$: 10 kurang dari 5, merupakan pernyataan yang salah, ditulis $\tau(\textit{q}) = F$.
    \end{enumerate}
  \end{contoh}
\end{frame}

\section{Pernyataan Majemuk}
\begin{frame}
\frametitle{Pernyataan Majemuk}
  \begin{definisi}[Pernyataan Majemuk]
   Suatu pernyataan yang didalamnya terdiri dari beberapa pernyataan, antara pernyataan yang satu dengan yang lainnya dihubungkan dengan kata sambung. Kata sambung ini berupa kata sambung logika.
  \end{definisi}
  
  \begin{block}{Kata Sambung Logika}
  	\begin{tabular}{ | l | l | l | }
  		\hline
  		Kata Sambung & Simbol & Nama \\ \hline
  		Dan (\textit{And}) & $\wedge$ & Konjungsi \\ \hline
  		Atau (\textit{or}) & $\vee$ & Disjungsi \\ \hline
  		Tidak (\textit{not}) & $\neg$ & Negasi \\ \hline
  		Jika ... maka ...(\textit{if ... then ...}) & $\rightarrow$ & Implikasi \\ \hline
  		jika dan hanya jika (\textit{if and only if}) & $\leftrightarrow$ & Biimplikasi \\ 
  		\hline
  	\end{tabular}
  \end{block}
  
  
\end{frame}

\subsection{Nilai Kebenaran Konjungsi}
\begin{frame}[shrink]
\frametitle{Nilai Kebenaran Konjungsi}
   \begin{definisi}[Konjungsi]
   	Adalah pernyataan yang dibentuk dari dua pernyataan \textit{p} dan \textit{q} dengan kata hubung "dan".   	
   \end{definisi}
    \begin{block}{Nilai Kebenaran Konjungsi}
    	\begin{tabular}{ | l | l | l | }
    		\hline
    		\textit{p} & \textit{q} & $\textit{p}\wedge\textit{q}$ \\ \hline
    		T & T & T \\ \hline
    		T & F & F \\ \hline
    		F & T & F \\ \hline
    		F & F & F \\     		 
    		\hline
    	\end{tabular}
    \end{block}
    \begin{contoh}
        Pernyataan : 13 bilangan prima dan $13^2$ = 169 \\
        Disimbolkan :\\
        \textit{p} : 13 bilangan prima ($\tau(\textit{p})$ =T)\\
        \textit{q} : $13^2$ = 169 ($\tau(\textit{q})$ =T)\\
        $\textit{p}\wedge\textit{q}$ : 13 bilangan prima dan $13^2$ = 169 ($\tau(\textit{p}\wedge\textit{q})$ =T)
    \end{contoh}
    
\end{frame}

\subsection{Nilai Kebenaran Disjungsi}
\begin{frame}[shrink]
	\frametitle{Nilai Kebenaran Disjungsi}
	\begin{definisi}[Disjungsi]
		Adalah pernyataan yang dibentuk dari dua pernyataan \textit{p} dan \textit{q} dengan kata hubung "atau".   	
	\end{definisi}
	\begin{block}{Nilai Kebenaran Disjungsi}
		\begin{tabular}{ | l | l | l | }
			\hline
			\textit{p} & \textit{q} & $\textit{p}\vee\textit{q}$ \\ \hline
			T & T & T \\ \hline
			T & F & T \\ \hline
			F & T & T \\ \hline
			F & F & F \\     		 
			\hline
		\end{tabular}
	\end{block}
	
	\begin{contoh}
		Pernyataan : 6 adalah bilangan genap atau 13 adalah bilangan prima.
		 \\
		Disimbolkan :\\
		\textit{p} : 6 adalah bilangan genap ($\tau(\textit{p})$ =T)\\
		\textit{q} : 13 adalah bilangan prima ($\tau(\textit{q})$ =T)\\
		$\textit{p}\vee\textit{q}$ : \\6 adalah bilangan genap atau 13 adalah bilangan prima ($\tau(\textit{p}\vee\textit{q})$ =T)
	\end{contoh}	
	
\end{frame}

\subsection{Nilai Kebenaran Negasi}
\begin{frame}[shrink]
	\frametitle{Nilai Kebenaran Negasi}
	\begin{definisi}[Negasi]
		Adalah pernyataan yang menyangkal atau mengingkari pernyataan awal.		   	
	\end{definisi}
	\begin{block}{Nilai Kebenaran Negasi}
		\begin{tabular}{ | l | l | }
			\hline
			\textit{p} & $\neg\textit{p}$  \\ \hline
			T & F  \\ \hline			
			F & T  \\     		 
			\hline
		\end{tabular}
	\end{block}
	\begin{contoh}
		 \textit{p}  : 2 + 3 = 5 ($\tau(\textit{p})$ = T)\\
		 $\neg\textit{p}$ : 2 + 3 $\neq$ 5 ($\tau(\neg\textit{p})$ = F)		 
	\end{contoh}
	\begin{contoh}
	 q: Semua bilangan prima adalah ganjil ($\tau(\textit{q})$ = T) \\
	$\neg\textit{q}$: Tidak benar bahwa semua bilangan prima adalah ganjil ($\tau(\neg\textit{q})$ = F)\\
	atau\\
	$\neg\textit{q}$ 	: Ada bilangan prima yang tidak ganjil ($\tau(\neg\textit{q})$ = F)	 
	\end{contoh}
\end{frame}


\subsection{Nilai Kebenaran Implikasi}
\begin{frame}[shrink]
	\frametitle{Nilai Kebenaran Implikasi}
	\begin{definisi}[Implikasi]
		
		Adalah pernyataan yang dibentuk dari dua pernyataan \textit{p} dan \textit{q} dalam bentuk "jika \textit{p}, maka \textit{q}".     	
	\end{definisi}
	\begin{block}{Nilai Kebenaran Implikasi}
		\begin{tabular}{ | l | l | l | }
			\hline
			\textit{p} & \textit{q} & $\textit{p}\rightarrow\textit{q}$ \\ \hline
			T & T & T \\ \hline
			T & F & F \\ \hline
			F & T & T \\ \hline
			F & F & T \\     		 
			\hline
		\end{tabular}
	\end{block}
	\begin{contoh}
		Pernyataan : Jika 3 + 2 = 5, maka 5 adalah bilangan prima
		\\
		Disimbolkan :\\
		\textit{p} : 3 + 2 = 5 ($\tau(\textit{p})$ =T)\\
		\textit{q} : 5 adalah bilangan prima ($\tau(\textit{q})$ =T)\\
		$\textit{p}\rightarrow\textit{q}$ : \\Jika 3 + 2 = 5, maka 5 adalah bilangan prima ($\tau(\textit{p}\rightarrow\textit{q})$ =T)	
	\end{contoh}
\end{frame}



\subsection{Nilai Kebenaran Biimplikasi}
\begin{frame}[shrink]
	\frametitle{Nilai Kebenaran Biimplikasi}
	\begin{definisi}[Biimplikasi]
			Adalah pernyataan yang dibentuk dari dua pernyataan \textit{p} dan \textit{q} dalam bentuk "\textit{p} jika dan hanya jika \textit{q}". 	
	\end{definisi}
	\begin{block}{Nilai Kebenaran Biimplikasi}
		\begin{tabular}{ | l | l | l | }
			\hline
			\textit{p} & \textit{q} & $\textit{p}\leftrightarrow\textit{q}$ \\ \hline
			T & T & T \\ \hline
			T & F & F \\ \hline
			F & T & F \\ \hline
			F & F & T \\     		 
			\hline
		\end{tabular}
	\end{block}
	\begin{contoh}
			Pernyataan : $16^{1/2} = 4$ jika dan hanya jika $\log^{16}{4} = \frac{1}{2}$\\
			Disimbolkan :\\
			\textit{p} : $16^{1/2} = 4$ ($\tau(\textit{p})$ =T)\\
			\textit{q} : $\log^{16}{4} = \frac{1}{2}$ ($\tau(\textit{q})$ =T)\\
			$\textit{p}\leftrightarrow\textit{q}$ : \\$16^{1/2} = 4$ jika dan hanya jika $\log^{16}{4} = \frac{1}{2}$ ($\tau(\textit{p}\leftrightarrow\textit{q})$ =T)	
	\end{contoh}
\end{frame}

\subsection{Varian dari Implikasi}
\begin{frame}
	\frametitle{Varian dari Implikasi}
  \small	
	\begin{block}{Konvers}
	$q \rightarrow p$, yang disebut konvers dari $\textit{p} \rightarrow \textit{q}$.
     \end{block}	
	\begin{block}{Invers}
	$\neg\textit{p}\rightarrow\neg\textit{q}$, yang disebut invers dari $\textit{p} \rightarrow \textit{q}$.\\
    \end{block}	
	\begin{block}{Kontraposisi}
	$\neg\textit{q} \rightarrow\neg\textit{p}$, yang disebut kontraposisi dari $\textit{p} \rightarrow \textit{q}$.\\
    \end{block}
	\begin{block}{Tabel Kebenaran Konvers,Invers dan Kontraposisi}	
		\begin{tabular}{ | l | l | l | l | l | l | l |l |}
			\hline
			\textit{p} & $\neg\textit{p}$ & \textit{q} & $\neg\textit{q}$ & $\textit{p}\rightarrow\textit{q}$ & $\textit{q}\rightarrow\textit{p}$ & $\neg\textit{p}\rightarrow\neg\textit{q}$ & $\neg\textit{q}\rightarrow\neg\textit{p}$\\ \hline
			T & F & T & F & T & T & T & T \\ \hline
			T & F & F & T & F & T & T & F \\ \hline
			F & T & T & F & T & F & F & T \\ \hline
			F & T & F & T & T & T & T & T \\     		 
			\hline
		\end{tabular}
   \end{block}	
	
\end{frame}

\subsection{Tautologi dan Kontradiksi}
\begin{frame}[shrink]
	\frametitle{Tautologi dan Kontradiksi}

	\begin{block}{Tautologi}
    Pernyataan majemuk disebut \textbf{tautologi} jika ia benar untuk semua kasus.    
    \end{block}	
    \begin{contoh}
      $p \vee \neg(p \wedge q)$ adalah sebuah tautologi\\
      \begin{tabular}{ | l | l | l | l | l |}
      	\hline
      	\textit{p} & \textit{q} &  $\textit{p}\wedge\textit{q}$ & $\neg(\textit{p}\wedge\textit{q})$ & $p \vee \neg(p \wedge q)$ \\ \hline
      	T &  T & T & F & T\\ \hline
      	T &  F & F & T & T \\ \hline
      	F &  T & F & T & T\\ \hline
      	F &  F & F & T & T\\     		 
      	\hline
      \end{tabular}
    \end{contoh}
    \begin{block}{Kontradiksi}
    	Pernyataan majemuk disebut \textbf{kontradiksi} jika ia salah untuk semua kasus
    \end{block}
    \begin{contoh}
    $(p \wedge q)\wedge\neg(p \vee q)$ adalah sebuah kontradiksi\\
    \begin{tabular}{ | l | l | l | l | l | l |}
    	\hline
    	\textit{p} & \textit{q} &  $\textit{p}\wedge\textit{q}$ & 
    	$\textit{p}\vee\textit{q}$ &
    	$\neg(\textit{p}\vee\textit{q})$ & $(p \wedge q)\wedge\neg(p \vee q)$ \\ \hline
    	T &  T & T & T & F & F\\ \hline
    	T &  F & F & T & F & F \\ \hline
    	F &  T & F & T & F & F\\ \hline
    	F &  F & F & F & T & F\\     		 
    	\hline
    \end{tabular}	
    \end{contoh}
\end{frame}	


\section{Negasi Pernyataan Majemuk}
\subsection{Negasi Konjungsi}
\begin{frame}
	\frametitle{Negasi Konjungsi}
    	\begin{definisi}[Negasi Konjungsi]
    		Negasi dari pernyataan $\textit{p}\wedge\textit{q}$  adalah $\neg\textit{p}\vee\neg\textit{q}$    		 	
    	\end{definisi}
    	\begin{contoh}
    		\textit{p} : saya suka apel \\
    		\textit{q} : saya tidak suka wortel \\
    		$\textit{p}\wedge\textit{q}$ : saya suka apel dan tidak suka wortel\\
    		$\neg(\textit{p}\wedge\textit{q})$ : saya tidak suka apel atau saya suka wortel	\\
    		
    		\begin{tabular}{ | l | l | l | l | l | l | l |}
    			\hline
    			\textit{p} & $\neg\textit{p}$ & \textit{q} & $\neg\textit{q}$ & $\textit{p}\wedge\textit{q}$ & $\neg(\textit{p}\wedge\textit{q})$ & $\neg\textit{p}\vee\neg\textit{q}$ \\ \hline
    			T & F & T & F & T & F & F\\ \hline
    			T & F & F & T & F & T & T \\ \hline
    			F & T & T & F & F & T & T\\ \hline
    			F & T & F & T & F & T & T\\     		 
    			\hline
    		\end{tabular}
    	\end{contoh}
\end{frame}

\subsection{Negasi Disjungsi}
\begin{frame}
	\frametitle{Negasi Disjungsi}
    \begin{definisi}[Negasi Disjungsi]
    	Negasi dari pernyataan $\textit{p}\vee\textit{q}$  adalah $\neg\textit{p}\wedge\neg\textit{q}$    		 	
    \end{definisi}
    \begin{contoh}
    	\textit{p} : Andi pergi ke supermarket \\
    	\textit{q} : Andi menonton di bioskop \\
    	$\textit{p}\vee\textit{q}$ : Andi pergi ke supermarket atau menonton di bioskop\\
    	$\neg(\textit{p}\vee\textit{q})$ : Andi tidak pergi ke supermarket dan tidak menonton di bioskop	\\
    	
    	\begin{tabular}{ | l | l | l | l | l | l | l |}
    		\hline
    		\textit{p} & $\neg\textit{p}$ & \textit{q} & $\neg\textit{q}$ & $\textit{p}\vee\textit{q}$ & $\neg(\textit{p}\vee\textit{q})$ & $\neg\textit{p}\wedge\neg\textit{q}$ \\ \hline
    		T & F & T & F & T & F & F\\ \hline
    		T & F & F & T & T & F & F \\ \hline
    		F & T & T & F & T & F & F\\ \hline
    		F & T & F & T & F & T & T\\     		 
    		\hline
    	\end{tabular}
    \end{contoh}
\end{frame}

\subsection{Negasi Implikasi}
\begin{frame}
	\frametitle{Negasi Implikasi}
 \begin{definisi}[Negasi Implikasi]
 	Negasi dari pernyataan $\textit{p}\rightarrow\textit{q}$  adalah $\textit{p}\wedge\neg\textit{q}$    		 	
 \end{definisi}
 \begin{contoh}
 	\textit{p} : Nico belajar dengan giat \\
 	\textit{q} : Nico naik kelas \\
 	$\textit{p}\rightarrow\textit{q}$ : Jika nico belajar dengan giat maka nico naik kelas\\
 	$\neg(\textit{p}\rightarrow\textit{q})$ : Nico belajar dengan giat dan ternyata nico tidak naik kelas.	\\
 	
 	\begin{tabular}{ | l | l | l | l | l | l | l |}
 		\hline
 		\textit{p} & $\neg\textit{p}$ & \textit{q} & $\neg\textit{q}$ & $\textit{p}\rightarrow\textit{q}$ & $\neg(\textit{p}\rightarrow\textit{q})$ & $\textit{p}\wedge\neg\textit{q}$ \\ \hline
 		T & F & T & F & T & F & F\\ \hline
 		T & F & F & T & F & T & T \\ \hline
 		F & T & T & F & T & F & F\\ \hline
 		F & T & F & T & T & F & F\\     		 
 		\hline
 	\end{tabular}
 \end{contoh}
\end{frame}

\subsection{Negasi Biimplikasi}
\begin{frame}[shrink]
	\frametitle{Negasi Biimplikasi}
\begin{definisi}[Negasi Biimplikasi]
	Negasi dari pernyataan $\textit{p}\leftrightarrow\textit{q}$  adalah $(\textit{p}\wedge\neg\textit{q})\vee(\textit{q}\wedge\neg\textit{p})$    		 	
\end{definisi}
\begin{contoh}
	\textit{p} : Ulangan dibatalkan \\
	\textit{q} : Diadakan kerja bakti \\
	$\textit{p}\leftrightarrow\textit{q}$ : Ulangan dibatalkan jika dan hanya jika diadakan kerja bakti \\
	$\neg(\textit{p}\leftrightarrow\textit{q})$ : Ulangan dibatalkan dan tidak diadakan kerja bakti atau diadakan kerja bakti dan ulangan tidak dibatalkan.	\\
 \tiny{	
	\begin{tabular}{ | l | l | l | l | l | l | l |l |l |}
		\hline
		\textit{p} & $\neg\textit{p}$ & \textit{q} & $\neg\textit{q}$ & $\textit{p}\leftrightarrow\textit{q}$ & $\neg(\textit{p}\leftrightarrow\textit{q})$ & $\textit{p}\wedge\neg\textit{q}$ & $\textit{q}\wedge\neg\textit{p}$ & $(\textit{p}\wedge\neg\textit{q})\vee(\textit{q}\wedge\neg\textit{p})$\\ \hline
		T & F & T & F & T & F & F & F & F\\ \hline
		T & F & F & T & F & T & T & F & T\\ \hline
		F & T & T & F & F & T & F & T & T\\ \hline
		F & T & F & T & T & F & F & F & F\\     		 
		\hline
	\end{tabular}}
\end{contoh}
\end{frame}

\section{Hirarki Operator Logika}
\begin{frame}
	\frametitle{Hirarki Operator Logika}
  \begin{definisi}[Hirarki Operator Logika]
  	Urutan pengerjaan dari operator logika.
  \end{definisi}
  \begin{block}{Tabel Hirarki Operator Logika }
  \begin{tabular}{ | l | l | l |}
  	\hline
  	Hierarki Ke & Simbol Operator & Nama Operator \\ \hline
  	1 & $\neg$ & Negasi \\ \hline
  	2 & $\wedge$ & Konjungsi \\ \hline
  	3 & $\vee$ & Disjungsi \\ \hline
  	4 & $\rightarrow$ & Implikasi \\ \hline
  	5 & $\leftrightarrow$ & Biimplikasi \\     		 
  	\hline
  \end{tabular}
  \end{block}
\end{frame}

\section{Latihan Soal}
\begin{frame}
	\frametitle{Latihan Soal}
	\textbf{SOAL 1}\\	
	Gunakan konstanta proposisional berikut:\\	 
	\begin{tabular}{l c l}
      \textit{p} & = & Bowo kaya raya\\
      \textit{q} & = & Bowo hidup bahagia
    \end{tabular}\\	
	Selanjutnya, ubah pernyataan-pernyataan berikut ini menjadi bentuk logika:\\	
	\begin{enumerate}
		\item Bowo tidak kaya.
		\item Bowo kaya raya dan hidup bahagia.
		\item Bowo kaya raya atau tidak hidup bahagia.
		\item Jika Bowo kaya raya, maka ia hidup bahagia.
		\item Bowo hidup bahagia jika dan hanya jika ia kaya raya.
	\end{enumerate}		
\end{frame}

\begin{frame}
	\frametitle{Latihan Soal}
	\textbf{SOAL 2}\\	
	Misalkan \textit{p},\textit{q} dan \textit{r} adalah variabel proposisional :\\	 
	\begin{tabular}{l c l}
		\textit{p} & = & Anda sakit flu\\
		\textit{q} & = & Anda Ujian\\
		\textit{r} & = & Anda Lulus
	\end{tabular}\\	
	Ubahlah ekspresi berikut menjadi pernyataan dalam bahasa indonesia : \\	
	\begin{enumerate}
		\item $p \rightarrow \neg q$
		\item $q \rightarrow \neg r$
		\item $\neg q \rightarrow r$
		\item $(p \wedge q) \rightarrow r$
		\item $(p \rightarrow \neg r) \vee (q \rightarrow \neg r)$
		\item $(p \wedge q) \vee (\neg q \wedge r)$
	\end{enumerate}		
\end{frame}

\begin{frame}
	\frametitle{Latihan Soal}
	\textbf{SOAL 3}\\	
	Ubahlah pernyataan-pernyataan berikut menjadi bentuk logika :\\	 	
	\begin{enumerate}
		\item Jika Bowo berada di Malioboro, maka Dewi juga ada di Malioboro.
		\item Pintu rumah Dewi berwarna merah atau coklat.
		\item Berita itu tidak menyenangkan.
		\item Bowo akan datang jika ia mempunyai kesempatan.
		\item Jika Dewi rajin kuliah, maka ia pasti pandai.		
	\end{enumerate}		
\end{frame}

\begin{frame}
	\frametitle{Latihan Soal}
	\textbf{SOAL 4}\\	
	Dengan menggunakan tabel kebenaran, jawab pertanyaan-pertanyaan berikut :\\	 	
	\begin{enumerate}
		\item Apakah nilai kebenaran dari $p \wedge p $ ?
		\item Apakah nilai kebenaran dari $p \vee p $ ?
		\item Apakah nilai kebenaran dari $(p \wedge \neg p) $ dan $(p \vee \neg p)$ ?
		\item Apakah $(p \rightarrow q)$ mempunyai nilai kebenaran yang sama dengan $(q \rightarrow p)$ ?
		\item Apakah $(p \rightarrow q) \rightarrow r$ mempunyai nilai kebenaran yang sama dengan $p \rightarrow (q \rightarrow r)$ ?		
	\end{enumerate}		
\end{frame}

\begin{frame}
	\frametitle{Latihan Soal}
	\textbf{SOAL 5}\\	
	Buatlah tabel kebenaran untuk ekspresi-ekspresi logika berikut :\\	 	
	\begin{enumerate}
		\item $\neg (\neg p \wedge \neg q)$
		\item $p \wedge (p \vee q)$
		\item $((\neg p \wedge (\neg q \wedge r))\vee (q \wedge r))\vee (p \wedge r)$
		\item $(p \wedge q)\vee (((\neg p \wedge q) \rightarrow p)\wedge \neg q)$
		\item $(p \rightarrow q)\leftrightarrow (\neg q \rightarrow \neg p)$		
		\item $p \wedge ((r \vee q) \leftrightarrow \neg r)$	
		\item $\neg((p \wedge q)\rightarrow \neg r)\vee p$	
	\end{enumerate}		
\end{frame}

\begin{frame}
	\frametitle{Latihan Soal}
	\textbf{SOAL 6}\\	
	Ubahlah pernyataan-pernyataan berikut menjadi ekspresi logika berupa proposisi majemuk : \\	 	
	\begin{enumerate}
		\item Jika tikus itu waspada dan bergerak cepat, maka kucing atau anjing itu tidak mampu menangkapnya.
		\item Jika saya tidak keliru, Dewi sudah diwisuda dan pacarnya atau orangtuanya berada disampingnya.
		\item Bowo membeli saham dan membeli properti untuk investasinya, atau dia dapat menanamkan uang di deposito bank dan menerima bunga uang.			
	\end{enumerate}		
\end{frame}

\begin{frame}
	\frametitle{Latihan Soal}
	\textbf{SOAL 7}\\	
    Masukan tanda kurung biasa ke dalam ekspresi logika berikut ini sehingga tidak terjadi ambiguitas : \\	 	
	\begin{enumerate}
		\item $p \wedge q \wedge r \rightarrow s$
		\item $p \vee q \vee r \leftrightarrow \neg s$
		\item $\neg p \wedge q \rightarrow \neg r \vee s$	
		\item $p \rightarrow q \leftrightarrow \neg r \rightarrow \neg s$
		\item $p \vee q \wedge r \rightarrow p \wedge q \vee r$		
	\end{enumerate}		
\end{frame}

\begin{frame}
	\frametitle{Latihan Soal}
	\textbf{SOAL 8}\\	
     Jika $\tau(p)=T$ dan $\tau(q)=T$,sedangkan $\tau(r)=F$ dan $\tau(s)=F$, carilah nilai kebenaran dari ekspresi-ekspresi logika berikut  : \\	 	
	\begin{enumerate}
		\item $p \wedge (q \vee r)$
		\item $(p \vee q) \wedge r$
		\item $((p\vee q)\wedge r) \vee \neg((p \vee q)\wedge (q\vee s))$	
		\item $(\neg(p \wedge q)\vee \neg r) \vee (((\neg p \wedge q)\vee \neg s) \wedge r)$
		\item $(p \leftrightarrow r) \wedge (\neg q \rightarrow s)$		
	\end{enumerate}		
\end{frame}

\begin{frame}
	\frametitle{Latihan Soal}
	\textbf{SOAL 9}\\	
	Jika $\tau(p)=F$ dan $\tau(q)=F$,sedangkan $\tau(r)=T$ dan $\tau(s)=T$, carilah nilai kebenaran dari ekspresi-ekspresi logika berikut  : \\	 	
	\begin{enumerate}
		\item $\neg(p \leftrightarrow q)\wedge (\neg r \rightarrow s)$
		\item $(\neg p \wedge (q \vee \neg r)\vee (q \leftrightarrow \neg p)) \leftrightarrow (s \wedge r)$
		\item $(p \vee (q \rightarrow (r \wedge \neg p))) \leftrightarrow \neg (q \vee \neg s)$		
	\end{enumerate}		
\end{frame}

\section{Ekuivalen Logis}
\begin{frame}
   \frametitle{Ekuivalen Logis}
   \begin{definisi}[Ekuivalen Logis]
   	 Preposisi \textit{p} dan \textit{q} disebut ekuivalen jika $\tau(p)=\tau(q)$, dengan simbol $p \equiv q$
   \end{definisi}
   
   \begin{contoh}
   	$\neg(p \wedge q) \equiv \neg q \vee \neg p$\\
   	Dibuktikan dengan tabel kebenaran\\
   	\begin{tabular}{ | l | l | l | l | l | l | l |}
   		\hline
   		\textit{p} & $\neg p$& \textit{q} &$\neg q$& $p \wedge q$  & $\neg(q \wedge p)$& $\neg q \vee \neg p$\\ \hline
   		T & F & T & F & T & F & F \\ \hline
   		T & F & F & T & F & T & T  \\ \hline
   		F & T & T & F & F & T & T\\ \hline
   		F & T & F & T & F & T & T \\     		 
   		\hline
   	\end{tabular}
   \end{contoh}	
   
\end{frame}

\begin{frame}
	\frametitle{Ekuivalen Logis}	
	\begin{contoh}
		$p \rightarrow q \equiv \neg p \vee q$\\
		Dibuktikan dengan tabel kebenaran\\
		\begin{tabular}{ | l | l | l | l |l |}
			\hline
			\textit{p} & $\neg p$ &  \textit{q} & $p \rightarrow q$  & $\neg p \vee q$\\ \hline
			T & F & T & T & T  \\ \hline
			T & F & F & F & F   \\ \hline
			F & T & T & T & T \\ \hline
			F & T & F & T & T  \\     		 
			\hline
		\end{tabular}
	\end{contoh}
\end{frame}	
	
\begin{frame}
	\frametitle{Ekuivalen Logis}	
	\begin{contoh}
		$p \leftrightarrow q \equiv (p \rightarrow q)\wedge (q \rightarrow p)$\\
		Dibuktikan dengan tabel kebenaran\\
		\begin{tabular}{ | l | l | l | l |l |l |}
			\hline
			\textit{p} & \textit{q} & $p \rightarrow q$  & $q \rightarrow p$ & $(p \rightarrow q)\wedge (q \rightarrow p)$ & $p \leftrightarrow q$\\ \hline
			T & T & T & T & T & T  \\ \hline
			T & F & F & T & F & F  \\ \hline
			F & T & T & F & F & F\\ \hline
			F & F & T & T & T & T \\     		 
			\hline
		\end{tabular}
	\end{contoh}	
\end{frame}

\begin{frame}
	\frametitle{Ekuivalen Logis}	
	\begin{contoh}
		$p \leftrightarrow q \equiv (p \wedge q)\vee(\neg p \wedge \neg q)$\\
		Dibuktikan dengan tabel kebenaran\\
		\begin{tabular}{ | l | l | l | l | l | l |l |l |}
			\hline
			\textit{p} & $\neg p$ &\textit{q}& $\neg q$ & $p \wedge q$  & $\neg q \wedge \neg p$ & $(p \wedge q)\vee(\neg p \wedge \neg q)$ & $p \leftrightarrow q$\\ \hline
			T & F & T & F & T & F & T & T  \\ \hline
			T & F & F & T & F & F & F & F  \\ \hline
			F & T & T & F & F & F & F & F\\ \hline
			F & T & F & T & F & T & T & T \\     		 
			\hline
		\end{tabular}
	\end{contoh}	
\end{frame}

\section{Latihan Soal}
\begin{frame}
	\frametitle{Latihan Soal}
	\textbf{SOAL 10}\\
	Buktikan bahwa ekspresi-ekspresi logika berikut ini ekuivalen dengan menggunakan tabel kebenaran :\\
	\begin{enumerate}
		\item $ \neg p \leftrightarrow q  \equiv (\neg p \vee q)\wedge (\neg q \vee p) $
		\item $q \leftrightarrow (\neg p \rightarrow q) \equiv T$
		\item $(p \vee \neg q)\rightarrow r \equiv (\neg p \wedge q)\vee r$	
		\item $p \rightarrow (q \rightarrow r) \equiv (p \rightarrow r)\rightarrow r$
		\item $p \rightarrow q \equiv \neg (p \wedge \neg q) $
		\item $\neg (\neg(p \wedge q)\vee q) \equiv F$
		\item $ ((p \wedge (q\rightarrow r))\wedge(p \rightarrow(q \rightarrow \neg r)))\rightarrow p \equiv T$	
	\end{enumerate}	
\end{frame}


\section{Hukum-hukum Logika}
\begin{frame}
	\frametitle{Hukum-hukum Logika}
	\begin{block}{Hukum-hukum Logika}
	\small	
	   \begin{tabular}{ | l | l |}
	   	\hline
	   	 Hukum Identitas  & Hukum Null Dominasi  \\ 
	   	 $p\vee F \equiv p$  & $p\wedge F \equiv F$  \\  
	   	 $p\wedge T \equiv p$  &  $p\vee T \equiv T$ \\ \hline
	   	  Hukum Negasi & Hukum Idempoten \\	
	   	  $p\vee \neg p \equiv T$  & $p\vee p \equiv p$  \\  
	   	  $p\wedge \neg q \equiv F$  &  $p\wedge p \equiv p$ \\ \hline
	   	  Hukum Involusi(Negasi Ganda) & Hukum Penyerapan (Absorpsi) \\	
	   	  $\neg(\neg p) \equiv p$  & $p \vee (p \wedge q) \equiv p$  \\  
	   	                                 &  $p \wedge (p \vee q) \equiv p$ \\ \hline
	   	  Hukum Komutatif & Hukum Asosiatif \\	
	   	  $p \vee q \equiv q \vee p$  & $ p \vee (q \vee r) \equiv (p \vee q) \vee r$  \\  
	   	  $p \wedge q \equiv q \wedge p$  &  $ p \wedge (q \wedge r) \equiv (p \wedge q) \wedge r$ \\ \hline	
	   	 Hukum Distributif & Hukum De Morgan \\	
	   	 $p \vee (q \wedge r) \equiv (p \vee q)\wedge (p \vee r) $  & $ \neg(p \wedge q) \equiv \neg p \vee \neg q$  \\  
	   	 $p \wedge (q \vee r) \equiv (p \wedge q)\vee (p \wedge r)$  &  $ \neg(p \vee q) \equiv \neg p \wedge \neg q$ \\
	   	\hline
	   \end{tabular}
	\end{block}
\end{frame}

\begin{frame}
	\frametitle{Hukum-hukum Logika}
	\begin{block}{Hukum-hukum Logika}
		\small	
		\begin{tabular}{ | l | l |}
			\hline
			Hukum Penyerapan (Absorpsi)  &   \\ 
			$p \vee (\neg p \wedge q) \equiv (p \vee q)$  & $(p \wedge q)\vee(p \wedge \neg q)\equiv p$  \\  
			$p \wedge (\neg p \vee q) \equiv (p \wedge q)$  &   \\ \hline
				
			$p \rightarrow q \equiv \neg p \vee q$  & $p \leftrightarrow q \equiv (p \wedge q) \vee (\neg p \wedge \neg q)$  \\  
			$p \rightarrow q \equiv \neg (p \wedge \neg q)$  &  $p \leftrightarrow q \equiv (p \rightarrow q) \wedge (q \rightarrow p)$ \\ \hline
				
			$(p \wedge q)\vee (p \wedge \neg q) \equiv p$  & $(p \wedge q)\vee (\neg p \wedge  q) \equiv q$  \\  
			$(p \vee q)\wedge (p \vee \neg q) \equiv p$& $(p \vee q)\wedge (\neg p \vee  q) \equiv q$ \\ \hline		
			
		\end{tabular}
	\end{block}
\end{frame}


\section{Penyederhanaan Ekspresi Logika}
\begin{frame}
	\frametitle{Penyederhanaan Ekspresi Logika}
	Penyederhana ekspresi logika di lakukan dengan menggunakan hukum logika.\\
	\begin{contoh}
		Sederhanakan $(p \wedge \neg q)\vee(p \wedge q \wedge r)$ :\\
		$(p \wedge \neg q)\vee(p \wedge q \wedge r)$\\		
		\begin{tabular}{ l l  l  }						
			$\equiv$ & $(p \wedge \neg q)\vee(p \wedge (q \wedge r))$ & Tambah Kurung   \\ 
			$\equiv$ & $p \wedge (\neg q \vee(q \wedge r))$ & Distributif   \\ 
			$\equiv$ & $p \wedge((\neg q \vee q)\wedge(\neg p \vee r)$ & Distributif \\ 
			$\equiv$ & $p \wedge(T \wedge (\neg p \vee r))$ & Tautologi  \\ 
			$\equiv$ & $p \wedge(\neg p \vee r)$ & Identitas   			
		\end{tabular}
	\end{contoh}	
\end{frame}

\begin{frame}
	\frametitle{Penyederhanaan Ekspresi Logika}
	Penyederhana juga dapat digunakan untuk membuktikan ekuivalensi logis.\\
	\begin{contoh}
		Buktikan : $(p \rightarrow q)\wedge(q \rightarrow p)\equiv(p \wedge q)\vee(\neg p \wedge \neg q)$ \\
		Pembuktian :\\
		$(p \rightarrow q)\wedge(q \rightarrow p)$\\		
		\begin{tabular}{ l l  l  }						
			$\equiv$ & $(\neg p \vee q)\wedge(\neg q \vee p)$ & $p \rightarrow q \equiv \neg p \vee q$   \\ 
			$\equiv$ & $(q \vee \neg p)\wedge(p \vee \neg q)$ & Komutatif   \\ 
			$\equiv$ & $(p \vee \neg q)\wedge(q \vee \neg p)$ & Komutatif \\ 
			$\equiv$ & $((p \vee \neg q)\wedge q)\vee((p \vee \neg q)\wedge \neg p)$ & Distributif  \\ 
			$\equiv$ & $((p \wedge q)\vee(\neg q \wedge q))\vee((p \wedge \neg p)\vee(\neg q \wedge \neg p))$ & Distributif  \\ 
			$\equiv$ & $((p \wedge q)\vee F)\vee(F\vee(\neg q \wedge \neg p))$ & Negasi   \\ 
			$\equiv$ & $(p \wedge q)\vee(\neg q \wedge \neg p)$ & Identitas   \\ 
			$\equiv$ & $(p \wedge q)\vee(\neg p \wedge \neg q)$ & Komutatif  			
		\end{tabular}
	\end{contoh}	
\end{frame}

\begin{frame}[shrink]
	\frametitle{Penyederhanaan Ekspresi Logika}
	Penyederhana juga dapat digunakan untuk memeriksa tautologi.\\
	\begin{contoh}
		 $((p \wedge(q \rightarrow r))\wedge(p \rightarrow(q \rightarrow \neg r)))\rightarrow p$ \\				
		\begin{tabular}{ l l  l  }						
			$\equiv$ & $\neg((p \wedge (\neg q \vee r))\wedge(\neg p \vee(\neg q \vee \neg r)))\vee p$ & $p \rightarrow q \equiv \neg p \vee q$   \\ 
			$\equiv$ & $(\neg(p \wedge (\neg q \vee r))\vee \neg(\neg p \vee(\neg q \vee \neg r)))\vee p$ & De Morgan   \\  
			$\equiv$ & $((\neg p \vee \neg(\neg q \vee r))\vee (\neg\neg p \wedge \neg(\neg q \vee \neg r)))\vee p$ & De Morgan   \\  
			$\equiv$ & $((\neg p \vee (\neg\neg q \wedge \neg r))\vee (\neg\neg p \wedge (\neg\neg q \wedge \neg\neg r)))\vee p$ & De Morgan   \\
			$\equiv$ & $((\neg p \vee (q \wedge \neg r))\vee  (p \wedge  (q \wedge r)))\vee p$ & Negasi Ganda   \\  	
			$\equiv$ & $(\neg p \vee (q \wedge \neg r))\vee  (p \wedge  (q \wedge r))\vee p$ & Hapus Kurung   \\  
			$\equiv$ & $(\neg p \vee (q \wedge \neg r))\vee ((p \wedge  (q \wedge r))\vee p)$ & Tambah Kurung   \\ 
			$\equiv$ & $(\neg p \vee (q \wedge \neg r))\vee (p \vee (p \wedge  (q \wedge r)))$ & Komutatif   \\ 
			$\equiv$ & $(\neg p \vee (q \wedge \neg r))\vee p $ & Absorption   \\ 
			$\equiv$ & $((\neg p \vee (q \wedge \neg r))\vee p) $ & Tambah Kurung   \\ 	
			$\equiv$ & $(p \vee(\neg p \vee (q \wedge \neg r))) $ & Komutatif   \\ 
			$\equiv$ & $(p \vee \neg p) $ & Absorpsi   \\
			$\equiv$ & $T$ & Tautologi   \\			
		\end{tabular}
	\end{contoh}	
\end{frame}

\section{Latihan Soal}
\begin{frame}
	\frametitle{Latihan Soal}
	\textbf{SOAL 11}\\
	Sederhanakan ekspresi-ekspresi logika berikut ini menjadi bentuk paling sederhana :\\
	\begin{enumerate}
		\item $ p \wedge (\neg p \rightarrow p) $
		\item $\neg (\neg p \wedge (q\vee \neg q))$
		\item $(p \rightarrow q)\rightarrow((p \rightarrow \neg q)\rightarrow \neg p)$	
		\item $(p \rightarrow(p \vee \neg r))\wedge \neg p \wedge q$
		\item $(\neg p \wedge \neg (q \rightarrow r))\wedge \neg (p \rightarrow (q \rightarrow \neg r)) $
		\item $(\neg p \rightarrow q)\rightarrow ((\neg p \rightarrow \neg q)\rightarrow p)$
		\item $(p \wedge (\neg p \vee q))\vee q \vee (p \wedge (p \vee q)) $
		\item $(p \wedge \neg q \wedge p \wedge \neg r)\vee (r \wedge p \wedge \neg r) $	
	\end{enumerate}	
\end{frame}


\begin{frame}
	\frametitle{Latihan Soal}
	\textbf{SOAL 12}\\
	Buktikan absorption laws berikut ini dengan penyederhanaan :\\
	\begin{enumerate}
		\item $ (p \wedge q)\vee(\neg p \wedge q) \equiv q $
		\item $ (p \vee q)\wedge(\neg p \vee q) \equiv q $
	\end{enumerate}	
	\textbf{SOAL 13}\\
	Sederhanakan ekspresi-ekspresi logika berikut ini menjadi bentuk paling sederhana :\\
	\begin{enumerate}
		\item $ \neg p \rightarrow \neg q $
		\item $(p \rightarrow q) \wedge (q \rightarrow r)$	
		\item $(p \rightarrow q)\leftrightarrow ((p \wedge q) \leftrightarrow p)$	
	\end{enumerate}
\end{frame}

\begin{frame}[shrink]
	\frametitle{Latihan Soal}
	\textbf{SOAL 14}\\
	Buktikan dua ekspresi logika berikut ekuivalen dengan penyederhanaan :\\
	\begin{enumerate}
		\item $ (p \wedge q)\vee (q \wedge r) \equiv q \wedge (p \vee r) $
		\item $ \neg(\neg(p \wedge q)\vee p) \equiv T$
		\item $\neg(\neg p \vee \neg(r \vee s))\equiv (p \wedge r)\vee (p \wedge s)$	
		\item $p \wedge (\neg p \vee q)\equiv p \wedge q$
		\item $\neg(\neg p \wedge \neg q)\equiv p \vee q$
		\item $(p \wedge q)\vee(p \wedge \neg q)\equiv p$
		\item $p \leftrightarrow q \equiv \neg(\neg(\neg p \vee q)\vee \neg(\neg q \vee p)) $
		\item $p \leftrightarrow q \equiv (\neg p \vee q) \wedge (\neg q \vee p) $
		\item $ \neg(\neg p \vee \neg(q \vee r))\equiv (p \wedge q) \vee (p \wedge r) $
		\item $(p \vee r)\wedge(q \vee s)\equiv (p \wedge q)\vee(p \wedge s)\vee(q \wedge r)\vee(r \wedge s) $			
	\end{enumerate}	
\end{frame}

\begin{frame}
	\frametitle{Latihan Soal}
	\textbf{SOAL 15}\\
	Dengan menggunakan penyederhanan, periksa apakah ekspresi logika berikut merupakan tautologi :\\
	\begin{enumerate}
		\item $p \wedge (\neg p \rightarrow p) $
		\item $(p \rightarrow q)\rightarrow((p \rightarrow \neg q)\rightarrow \neg p) $	
		\item $((p \rightarrow (q \vee \neg r))\wedge \neg p)\wedge q $
		\item $(\neg p \rightarrow q)\rightarrow((\neg p \rightarrow \neg q)\rightarrow p) $	
		\item $((p \wedge (q \rightarrow r))\wedge(p \rightarrow(q \rightarrow \neg r)))\rightarrow p $
		\item $((p \rightarrow q)\wedge(q \rightarrow r))\rightarrow((p \wedge q)\rightarrow r) $
		\item $p \rightarrow (q \rightarrow p) $
		\item $(q \rightarrow p) \rightarrow p $
		\item $\neg \neg p \rightarrow p$
		\item $(\neg p \rightarrow \neg q)\rightarrow(q \rightarrow p)$	
		\item $(p \rightarrow (q \rightarrow r))\rightarrow((p \rightarrow q)\rightarrow(p \rightarrow r))$			
	\end{enumerate}	
\end{frame}

		

\end{document}